\documentclass{article}
\usepackage[preprint,main]{neurips_2026}
\usepackage[T1]{fontenc}
\renewcommand{\ttdefault}{lmtt}
\usepackage[utf8]{inputenc}
\usepackage{booktabs,tabularx,amsmath}
\usepackage[colorlinks=true,linkcolor=black,citecolor=black,urlcolor=blue]{hyperref}
\title{Cross-Hospital Generalization of a Sepsis Early-Warning Classifier}
\author{Trisha Bajpai \qquad Basmala Metwally \qquad Jacob Loe\\Virginia Tech --- CS 4824 Milestone Report}
\date{}
\hypersetup{pdftitle={Cross-Hospital Generalization of a Sepsis Early-Warning Classifier},pdfauthor={Trisha Bajpai, Basmala Metwally, Jacob Loe},pdfsubject={CS 4824 Milestone Report}}
\begin{document}
\maketitle
\begin{abstract}
Across three seeds on PhysioNet 2019, logistic-regression utility averaged 0.2460 on Hospital A holdout and 0.0756 on B. It outperformed a partial SIRS comparator, but subgroup deficits limit reliance. An exploratory abstention rule withheld 3.54\% of B hours without demonstrating a utility gain.
\end{abstract}

\section{Question, thesis, and progress from the proposal}
The model would prioritize clinician review of ICU patients, not diagnose or treat. We compare within-hospital random evaluation with hospital transfer, motivated by record instability \citep{nestor2019}. Since the Challenge examined hospital generalization \citep{reyna2020}, this is replication and diagnosis.

\textbf{Current thesis.} A sepsis classifier may beat a simple rule yet lose utility across hospitals, so clinician-review reliance requires evidence on transfer, subgroups, calibration, and the consequences of abstention.

Following the proposal's fallback, we used regularized logistic regression instead of boosted trees for a reproducible, inspectable baseline within available compute. Boosted trees are not implemented. Jacob took over ingestion/documentation. A working pipeline, three-seed evaluation, and official scoring checks are complete; 18 tests pass. Code: \url{https://github.com/jacobloe1/sepsis-cross-hospital}.

\section{Implemented approach and evaluation design}
We used all 20,336 A and 20,000 B patients \citep{data2019}, preserving missing values and supplied labels, already advanced six hours. Seeded shuffles, stratified by ever-positive status, partition A patients 60/15/10/15\% into training/calibration/threshold-tuning/test sets (floor boundaries; 3,051 test patients). B supplies no development patients. A separate pooled A+B regime uses these fractions within hospital/sepsis strata (6,052 test patients); it exposes its separate model to B development data and does not isolate an evaluation-design effect.

The 108 features use current values, last observations, and missingness flags for 34 physiological variables, plus six demographic/time fields. Filling uses past/current information only. Training medians/means/scales impute and standardize inputs, clipped at $\pm10$. \textbf{Preprocessing and classifier fitting use a seeded reservoir capped at 50,000 training hours}; calibration/evaluation are uncapped. L2-regularized logistic regression uses 15 SGD epochs (L2=0.001); probability calibration uses separate patients.

The partial SIRS rule requires two criteria: HR$>90$; Temp$<36$ or $>38$; Resp$>20$ or PaCO$_2<32$; WBC$<4$ or $>12$. Unknown criteria do not count; bands are unavailable. Vitals expire after six hours and blood tests after 24; these assumptions await clinical review. We match the rule's label-negative-hour false-positive rate on tuning patients, conservatively handling ties, then freeze the model threshold. Equal evaluation burden is not guaranteed.

\newpage
\section{Preliminary findings, variation, and failed attempts}
\begin{table}[h]
\centering\small
\begin{tabular}{@{}lrrrr@{}}
\toprule
Evaluation/model & AUROC & Utility (mean $\pm$ SD) & Sensitivity & Coverage\\
\midrule
A / SIRS & 0.5897 & $0.0794\pm0.0074$ & 51.5\% & 100.0\%\\
A / logistic & 0.7257 & $0.2460\pm0.0239$ & 67.3\% & 100.0\%\\
B / SIRS & 0.6223 & $-0.0159\pm0.0000$ & 49.6\% & 100.0\%\\
B / logistic & 0.6951 & $0.0756\pm0.0174$ & 52.9\% & 100.0\%\\
B / logistic + gate & 0.6951 & $0.0741\pm0.0167$ & 52.4\% & 96.5\%\\
Pooled / logistic & 0.7424 & $0.2094\pm0.0162$ & 66.5\% & 100.0\%\\
\bottomrule
\end{tabular}

\caption{Means across seeds 7, 42, and 2026; sensitivity is hourly. \emph{Coverage} is the share of hours on which the model gives a prediction rather than abstaining. SD describes variation across fits/splits, \emph{not} a confidence interval. Gated AUROC uses unchanged latent risks on all hours. All B patients and their 761,995 hours are evaluated in every seed.}
\label{tab:results}
\end{table}

\paragraph{Transfer and uncertainty.}
Utility decreases from A holdout to B in every seed, even though AUROC changes less (Table~\ref{tab:results}). B logistic utility is .0713, .0608, and .0947. Improvements over SIRS have paired 95\% bootstrap intervals [.0618,.1129], [.0470,.1050], and [.0822,.1386], respectively. Each uses 1,000 whole-patient resamples with the fitted model/threshold held fixed, excluding training uncertainty. Across-seed SD instead combines changed patient partitions, reservoir samples, and SGD order. Three seeds are not independent hospital samples. Case mix, prevalence, and measurement practice remain unresolved explanations; this does not identify a causal splitting-strategy effect.

Our utility implementation agrees with the unmodified official NumPy scorer \citep{scorer} for observed, ideal-window, and no-action patient utilities and cohort normalization; maximum patient discrepancies are below $10^{-13}$. Warning timing is conditional on detection in the inferred onset $[-12,0]$h window (onset = first positive row + 6); it is not prospective lead-time evidence.

\paragraph{Abstention failure and exploratory revision.}
The abstention \emph{gate} withholds a prediction when recent data look too sparse to trust. The initial rule withheld hours whose raw missingness exceeded A-training's 95th percentile. That percentile was 1.0, so the strict inequality could never trigger. We retained this failure. \textbf{The revised analysis is exploratory because B had already been inspected.} It uses seven vital signs (HR, O2Sat, Temp, SBP, MAP, DBP, Resp) in the current and previous five hours. It withholds when none is recent, or their missing fraction exceeds the training 95th percentile among hours with at least one recent vital. Rule parameters are fixed before revised evaluation and learned without B in the A-only regime.

The revision activates but does not repair transfer: B coverage is 96.46\% in each seed, and utility changes versus ungated predictions are $-.00168$, $-.00069$, and $-.00217$. All three paired B intervals include zero: neither a benefit nor a definitive adverse B effect is established. Withheld hours remain in sensitivity/utility denominators as no alerts. Some A/pooled intervals exclude zero on the negative side. This freshness screen does not establish support for unusual laboratory patterns or a safe clinical fallback.

\paragraph{Calibration, subgroups, and current conclusion.}
Low overall ten-bin expected calibration error (ECE) hides sparse higher-risk failures: in original seed-42 B predictions, the .30--.40 bin averaged .340 predicted risk but .100 positive labels (521 hours). Unit-2 (patients flagged by the data's \texttt{Unit2} ICU-unit indicator) sensitivity is 7.96, 13.71, and 16.19 percentage points below overall B sensitivity; two fits exceed the proposal's ten-point limit on point estimates. Adequacy requires at least 30 septic and 30 never-septic patients, but subgroup intervals/power remain unresolved. B precision is 3.1--3.5\%, with mean 23.0 alert hours per 100 hours versus SIRS's 25.5. Label-negative alerts can include useful early warnings, so these are not direct clinical false-alarm estimates. Evidence supports a transfer-performance concern and an advantage over the partial rule, \emph{not} the proposal's full limited-reliance claim.

\newpage
\section{Draft class-survey case and comparison entry}
\paragraph{Case summary.}
This case concerns ICU sepsis warnings for clinician review. Evidence comes from public hourly records, with A-only development, held-out A evaluation, fixed B transfer, and a separate pooled random experiment. Logistic regression improves on a partial SIRS rule but loses utility on B across three seeds. A missingness gate first failed to activate; its revised exploratory version activates without demonstrating a utility gain. Unit-related sensitivity gaps and higher-risk calibration errors remain. The case illustrates why discrimination, operational burden, selective coverage, and subgroup evidence must be considered together. Human diagnosis and treatment remain outside the model's authority.

\begin{table}[h]
\centering\small
\begin{tabularx}{\linewidth}{@{}p{.21\linewidth}X@{}}
\toprule
Matrix field & Draft entry\\
\midrule
Domain / decision & Clinical ML; prioritize ICU patients for clinician sepsis review.\\
Evidence / split & PhysioNet 2019; A train/calibrate/tune/test; fixed B transfer; separate pooled random comparison.\\
Comparator / measures & Partial SIRS and logistic; utility, hourly sensitivity/PPV, calibration, burden, warning timing, coverage, and patient bootstrap.\\
Shift / failure & Hospital/measurement shift; inactive original abstention, no demonstrated gain from the revision, unit-2 deficits.\\
Lifecycle / conclusion & Clinician retains decisions; B-informed revisions are exploratory; limited reliance is not established.\\
\bottomrule
\end{tabularx}
\caption{Draft comparison-matrix entry, shown vertically for readability.}
\end{table}

\section{Completion plan and remaining obstacles}
Before the final report, we will proceed in the following order:
\begin{enumerate}
\item \textbf{Diagnose the unit-related failure.} Basmala will add patient-clustered subgroup intervals and review group sizes; Trisha will compare completeness and missing-unit indicators. Deliverable: a subgroup/measurement table separating observation from explanations. If groups remain imprecise, report uncertainty rather than certify the adequacy heuristic.
\item \textbf{Strengthen calibration and operational evaluation.} Jacob will add reliability diagnostics near the alert threshold and uncertainty for sparse bins; Trisha and Basmala will specify a cooldown/single-alert comparison and controlled vital-dropout stress test. Seek clinical input on the SIRS proxy and freshness assumptions; if unavailable, retain the proxy label and narrow its interpretation.
\item \textbf{Extend models after the baseline audit.} Jacob will attempt boosted trees only if compute permits, selecting within A development patients. Otherwise, retain logistic regression and analyze cap/convergence sensitivity. Preserve existing manifests/results; do not choose thresholds or a best seed to improve B. Profile memory or stream data if loading is limiting.
\item \textbf{Finalize the evidence record.} All three members will reconcile the thesis, matrix, contributions, citations, and limitations with results. Any B-informed revision remains exploratory. If new independent hospital evaluation is infeasible, report a bounded diagnostic result rather than claim untouched validation or clinical benefit.
\end{enumerate}

\section{Contributions}
\textbf{Jacob Loe} handled data acquisition, ingestion/documentation, classifier/calibration work, and original baseline execution, including ingestion responsibilities reassigned after the team change. \textbf{Trisha Bajpai} handled preprocessing and the SIRS-style baseline. \textbf{Basmala Metwally} handled evaluation, abstention, and reproducibility checks. These completed responsibilities were confirmed by Jacob. Initial implementation, troubleshooting, follow-up computations, and this draft used AI assistance; the team remains responsible for reviewing the analysis and final submission.

\label{lastcontent}

\clearpage
\begin{thebibliography}{9}
\bibitem[Nestor et~al.(2019)]{nestor2019}
B.~Nestor, M.~B.~A. McDermott, W.~Boag, G.~Berner, T.~Naumann, M.~C. Hughes, A.~Goldenberg, and M.~Ghassemi.
Feature robustness in non-stationary health records: Caveats to deployable model performance in common clinical machine learning tasks.
\emph{Proceedings of Machine Learning Research}, 106:381--405, 2019.
\url{https://proceedings.mlr.press/v106/nestor19a.html}.

\bibitem[PhysioNet(2019)]{scorer}
PhysioNet. Evaluation code for the PhysioNet/CinC Challenge 2019.
\url{https://github.com/physionetchallenges/evaluation-2019}, 2019.
The unmodified scorer and source hashes used here are retained in the project repository.

\bibitem[Reyna et~al.(2019)]{data2019}
M.~Reyna et~al.
Early Prediction of Sepsis from Clinical Data: The PhysioNet/Computing in Cardiology Challenge 2019, version 1.0.0.
PhysioNet, 2019.
\url{https://doi.org/10.13026/v64v-d857}.

\bibitem[Reyna et~al.(2020)]{reyna2020}
M.~A. Reyna, C.~S. Josef, R.~Jeter, S.~P. Shashikumar, M.~B. Westover, S.~Nemati, G.~D. Clifford, and A.~Sharma.
Early prediction of sepsis from clinical data: The PhysioNet/Computing in Cardiology Challenge 2019.
\emph{Critical Care Medicine}, 48(2):210--217, 2020.
\url{https://doi.org/10.1097/CCM.0000000000004145}.
\end{thebibliography}
\end{document}
